\documentclass[11pt,letterpaper]{article}

% Sourcing proposal, single column. Single column is deliberate: the central
% summary table has fifteen rows of prose, and the prior->posterior table has
% nine columns. Neither survives two columns legibly.

\usepackage[T1]{fontenc}
\usepackage{times}
\usepackage{amsmath,amssymb}
\usepackage{booktabs}
\usepackage{graphicx}
\usepackage[hidelinks]{hyperref}
\usepackage{microtype}
\usepackage{setspace}
\usepackage[letterpaper,top=1.0in,bottom=1.0in,left=1.0in,right=1.0in]{geometry}
\usepackage{caption}
\usepackage{xcolor}
\usepackage{listings}

\captionsetup{font=small,labelfont=bf}
\setlength{\parskip}{3pt}
\setstretch{1.04}
\newcommand{\code}[1]{\texttt{#1}}
\newcommand{\bs}{\allowbreak\hspace{0pt}}
\setlength{\tabcolsep}{5pt}

\lstdefinestyle{tiny}{basicstyle=\ttfamily\scriptsize,breaklines=true,
  columns=fullflexible,frame=single,rulecolor=\color{gray!40},xleftmargin=4pt,
  xrightmargin=4pt,aboveskip=4pt,belowskip=6pt}
\lstset{style=tiny}

\title{\textbf{Grounding the Decision Constants of a Self-Evaluating\\ Research System}\\[6pt]
\large A sourcing proposal: which values the literature determines,\\ which it
merely contains, and which it does not touch}
\author{
  \normalsize Prepared with \code{opencode}\\[3pt]
  \small Companion to \emph{Provisional Constants} (\code{paper.pdf}) and Figure~1
  (\code{flow.pdf})\\
  \small Audited system: \code{github.com/cloudcell/labloop}, commit \code{21d91f1}
}
\date{}

\begin{document}
\maketitle

\begin{abstract}
An audit of an agentic ML research system found fifteen numeric constants that
gate scientific decisions and carry no stated derivation. This document attempts,
for each, to locate literature that \emph{determines} the value. Four are
groundable in a specific source; two are settled by a theorem rather than by a
number; nine have no literature basis and should either be reclassified as
operational preferences or deleted. The largest single result is negative: a
constant producing a bounded score for a ratio whose denominator may be near
zero is \emph{invalid}, not merely conservative, by a theorem of 1987. The
second is quantitative: a session size of $n \approx 9$ seeds per arm is required
for 80\% directional accuracy at a specified smallest effect size of interest, a
figure the audited system's own benchmark independently confirms to within
$0.023$, against a shipped default of $n = 1$.
\end{abstract}

\section{Purpose}

The companion paper argues that a decision constant is not a parameter but a
claim, and that an unsourced claim should be labelled \code{PROVISIONAL} rather
than quietly published. This document performs the labelling. For every decision
constant identified in the audit, it asks a narrow question: \emph{is there a
retrievable source that fixes this value, and states why?} Where the answer is
no, that is recorded as a finding, not as a failure to search hard enough.

The result is stated up front so that the tables are not read as more favourable
than they are.

\begin{quote}
\textbf{Of fifteen decision constants audited, four can be grounded in a
specific citable source, two are settled by a mathematical theorem rather than
by a value, and nine have no literature basis at all --- and should be
reclassified as operational, or deleted.}
\end{quote}

A large fraction of the numbers that attracted audit attention are not
epistemic claims. They are timeouts, retention counts and bind addresses.
Presenting those to a literature review is a category error, and part of this
proposal's purpose is to say so explicitly rather than to manufacture citations
for them.

\section{What counts as grounding}

A constant is \code{SOURCED} only if a retrieved source states the value
\emph{and} the reason for it, such that a maintainer could substitute a
different value knowing what breaks. Three weaker outcomes are recorded
separately and must never be reported as grounding.

\begin{table}[h]
\centering
\caption{Status vocabulary. The third row is the failure mode this proposal is
most at risk of.}
\label{tab:status}
\begin{tabular}{@{}p{0.16\textwidth}p{0.76\textwidth}@{}}
\toprule
\textbf{Status} & \textbf{Meaning} \\
\midrule
\code{SOURCED} & A source names this value, or a procedure that yields it,
with justification. \\
\code{DERIVED} & No source gives the value, but a method does; the value
follows from inputs we supply. \\
\code{IN-RANGE} & A source tabulates a family of values and ours is one of
them, but the source does not single ours out. \\
\code{UNGROUNDED} & No source found. Delete, reclassify as operational, or
retain as \code{PROVISIONAL} with a stated reason. \\
\bottomrule
\end{tabular}
\end{table}

\code{IN-RANGE} is separated because it is how an audit of this kind deceives
itself: a paper is found that happens to contain the number $0.3$ somewhere in
a table of unrelated quantities, and the constant is declared grounded. The
distinction is maintained throughout \S\ref{sec:results}.

\section{C1 --- Calibration constants}

\subsection{The prior ceiling: \code{0.3}}
\label{sec:prior}

This is the best-anchored constant in the system. The National Academies' report
on reproducibility \cite{nap2019} tabulates the posterior probability that a
hypothesis is true as a function of the prior probability and the observed
$p$-value. Its Table~D-1 gives, for a one-tailed $p = 0.05$:

\begin{table}[h]
\centering
\caption{Posterior probability that the hypothesis is true, one-tailed
$p = 0.05$ (National Academies 2019, Table D-1). The bold column is the
audited system's prior ceiling.}
\label{tab:nap}
\footnotesize
\setlength{\tabcolsep}{4pt}
\begin{tabular}{@{}l*{8}{c}@{}}
\toprule
prior $P[H_1]$ & .01 & .05 & .10 & .20 & .25 & \textbf{.30} & .40 & .50 \\
\midrule
posterior & .038 & .169 & .301 & .492 & .563 & \textbf{.624} & .721 & .795 \\
\bottomrule
\end{tabular}
\end{table}

The report discusses $0.3$ explicitly as \emph{``its pre-experimental probability
of being true is about 1 in 3.''} That is a substantively defensible prior for a
scientific claim, and it maps cleanly onto a ceiling for a claim with no
evidence-bearing edge: an ungrounded claim is worth your prior, and one-in-three
is a documented choice rather than a guess.

\textbf{Proposal.} Retain $0.3$, but cite the source at the definition site,
make the assumed prior a deployable configuration value, and look the ceiling up
in the same table rather than hardcoding it. An operator who believes their prior
is $0.10$ should obtain a ceiling near $0.30$, not $0.3$.

\textbf{Caveat.} The report offers $0.3$ as one row among many, not as a
recommendation. The status is \code{IN-RANGE}, not \code{SOURCED}.

\subsection{Verdict confidence: \code{0.85}}

This is the constant the literature actively indicts rather than merely failing
to support. From the same table, at a one-tailed $p = 0.05$ the posterior
\emph{cannot exceed $0.795$} even at a $50\%$ prior. The audited system mints
$0.85$ on a verdict that may rest on a single trial, with no $p$-value computed
and no evidence count recorded. Reaching $0.85$ requires simultaneously a high
prior and a strong result; the system assumes the first by fiat and never measures
the second.

The report also states the principle the system violates \cite{nap2019}:

\begin{quote}\small
\emph{``It is clearly inappropriate to apply the same confidence to the results
of a study with a highly unexpected and surprising result as in a study in which
the results were a priori more plausible.''}
\end{quote}

The system applies an identical confidence to every accepted and every rejected
hypothesis, regardless of effect size, variance, or surprise --- which is the
specific thing the sentence forbids.

\textbf{Proposal.} Replace the verdict-to-float lookup with a computed posterior
using the deployment's declared prior and an effect-size test. For the common
case of prior $0.3$ at one-tailed $p = 0.05$, Table~D-1 gives \textbf{$0.624$}.
Where no $p$-value is available, the correct output is an enumeration, not a
float. A further check the system currently lacks is available: the Brier score
decomposes into \emph{calibration} and \emph{refinement} components
\cite{gneiting2007}, so the system could audit whether its confidences are
calibrated rather than assert that they are.

\subsection{Meta-decision confidence: \code{0.6}}

No literature supports $0.6$, and the provenance trace is worse than merely
unsourced. The value is traceable to a unit-test fixture whose assertion is that
the value \emph{must be rejected}:

\begin{lstlisting}
r = await call_tool(mcp, "record_finding", {
    "investigation_id": inv_id, "content": "too confident",
    "confidence": 0.6,
})
assert "prior ceiling" in r["error"]
\end{lstlisting}

\noindent
(\code{tests/zetesis/test\_20260916-1721Z--enforcement.py:109})

The production call site sits in \code{arete/tools/decisions.py:122} beside a
comment explaining that evidence edges are minted inline because the claims
server caps unevidenced claims at the prior ceiling. The mechanism is documented;
the magnitude is not.

\textbf{Proposal.} Delete. Replace with the same posterior computation as above,
keyed on the decision's own evidence count. Meta-decisions are required to cite
at least one evidence reference, so the inputs exist --- they are simply never
read.

\subsection{Importance and edge weight}

\code{edge.weight = 1} is defensible \emph{if} it is read as \emph{unweighted},
which is the honest interpretation: the claim graph does not attempt to rank
evidence. The difficulty is that a field named \emph{weight} carrying a constant
invites the reading that evidence is being weighted when it is not. Rename it to
declare the fact, or drop it.

\code{importance = 0.5} has no defensible reading. It is the midpoint of the unit
interval and it responds to nothing.

\textbf{Proposal.} \textbf{Delete \code{importance}.} No citation is needed to
delete a field that carries no information --- and that is a better outcome than a
spurious reference.

\section{C2 --- Discriminative constants}

\subsection{Zero-divisor handling: settled by theorem}

This is the one item where the literature supplies no number but \emph{proves the
current implementation wrong}. The proof has been available since 1954.

Fieller's construction of a confidence interval for a ratio of two means yields an
\emph{unbounded} interval whenever the denominator is not significantly different
from zero \cite{fieller1954} --- either the whole real line, or two disconnected
infinite rays. This is not a defect of the method; it is the correct answer, since
a ratio whose denominator is compatible with zero can be arbitrarily large in
either direction.

Gleser and Hwang \cite{gleser1987} then close the door on the alternative. A
retrieved restatement gives their result as a theorem:

\begin{quote}\small
\emph{``Any confidence interval which has finite length for almost every
observation has minimum zero coverage probability while minimizing over
$(\alpha_1,\alpha_2)$.''}
\end{quote}

Read plainly: \emph{one cannot manufacture a bounded, valid score for a ratio when
the denominator may be near zero.} Any method that always returns a finite number
in that situation has zero coverage. It is not a conservative estimate; it is an
invalid one. Returning $0.0$ is a decision that has been taken, not a limit that
has been reached. Independent reviews reach the same conclusion and note in
addition that the naive ratio estimator is biased away from the singular case, so
the problem is not confined to exact zero \cite{axioms2025,fiellerguide}.

Empirically this is what was observed. One campaign closed with a frozen score of
$0$ for \code{champion\bs \_mean = 0} against \code{challenger\bs \_mean = 0.5}: the
challenger beat the champion, the ratio is undefined, and the frozen score is the
worst value the scale admits. The same session produced a tournament, driven with
the identical input, that correctly refused to close.

\textbf{Proposal.}
\begin{enumerate}
\item The campaign closer must \emph{refuse} a zero divisor, matching the
tournament closer, which already does. This is a one-line correctness fix and the
single highest-value change identified by the audit.
\item Replace the point ratio with a Fieller confidence set and record the set's
boundedness. The correct output for a near-zero denominator is \emph{undefined}.
\item Refusing on exact equality to zero catches only the measure-zero case. Gate
on the Fieller condition instead: the denominator's confidence interval must
exclude zero. That condition covers the whole neighbourhood, which is where the
$37.8\%$ sign-error rate lives.
\end{enumerate}

\subsection{Session size $n$: derived from a citable procedure}
\label{sec:n}

This is the constant with a fully specified literature procedure, and the
procedure is unambiguous: declare a smallest effect size of interest, then compute
the $n$ that achieves a stated power against it \cite{lakens2018sesoi,albers2018}.
Two warnings from the same literature apply directly.

\emph{First, do not use the observed effect size from a pilot to set $n$.}
Albers and Lakens \cite{albers2018} demonstrate that pilot-derived effect sizes are
upwardly biased by follow-up bias, which systematically yields underpowered main
studies. The audited system's tournaments routinely run one seed and record what
they observe; that is the pilot-as-design error.

\emph{Second, relying on a benchmark is ``the weakest possible justification''} of a
smallest effect size of interest \cite{lakens2018sesoi}. The threshold must come
from the individual contract, not from a global default.

\subsubsection{Worked example, validated against the system's own data}

Using the parameters of the audit benchmark: base accuracy $\mu = 0.50$,
per-observation $\sigma = 0.05$, and a smallest effect size of interest
$\delta = 0.02$ in the challenger-minus-parent difference. On the log-ratio
scale $\delta \approx \log(0.52/0.50) = 0.03922$, and the delta-method standard
error is $\sigma\sqrt{2/n}/\mu$. Solving for the $n$ giving 80\% directional
accuracy ($z = 0.8416$):

\begin{equation}
n \;=\; 2\left(\frac{\sigma z}{\mu\,\delta}\right)^{2}
   \;=\; 2\left(\frac{0.05 \times 0.8416}{0.50 \times 0.03922}\right)^{2}
   \;=\; \mathbf{9.2} \ \ \text{seeds per arm}
\end{equation}

and $n \approx 21.4$ per arm for 90\%.

The normal approximation is not a hand-wave: the audited system's own Monte-Carlo
benchmark confirms it to within $0.023$ at every sample size tested
(Table~\ref{tab:valid}).

\begin{table}[h]
\centering
\caption{Closed-form prediction of sign error against measured Monte-Carlo sign
error, 400 replicates per cell, $\sigma = 0.05$, true effect $+0.02$.}
\label{tab:valid}
\begin{tabular}{@{}lccc@{}}
\toprule
\textbf{seeds / arm} & \textbf{predicted} & \textbf{measured} & \textbf{error} \\
\midrule
1  & 0.391 & 0.378 & $+0.013$ \\
5  & 0.268 & 0.245 & $+0.023$ \\
20 & 0.107 & 0.103 & $+0.004$ \\
\bottomrule
\end{tabular}
\end{table}

So the figure of roughly nine is not a formula transplanted from a textbook; it is
a prediction the audited system's own recorded data supports. The system ships
with a default of $n = 1$ --- a nine-fold deficit --- and records neither the
required nor the achieved $n$.

\textbf{Proposal.}
\begin{enumerate}
\item The promotion policy, currently inert, must require a smallest effect size of
interest and a target power, and the tournament must refuse to open without them.
This is what makes the field non-inert.
\item Store $n_{\text{requested}}$, $n_{\text{achieved}}$, the effect size of
interest, and the Type~S / Type~M error \cite{gelman2014} on the frozen score row
itself. A score is uninterpretable without them.
\end{enumerate}

\subsection{Promotion threshold: an odds ratio, not a metric ratio}

The current decision is \code{score > 1} --- an unweighted ratio with no threshold
search. The Bayes-factor literature supplies a graded evidence scale better suited
to the task, with named rungs a contract may select from
\cite{kass1995}:

\begin{table}[h]
\centering
\caption{Evidence rungs on the $2\ln BF$ scale (Kass \& Raftery 1995).}
\label{tab:bf}
\begin{tabular}{@{}lccccc@{}}
\toprule
$2\ln BF$ & 0--2 & 2--6 & 6--10 & 10--150 & $>150$ \\
\midrule
reading & not worth more & & & & \\
         & than a bare & positive & strong & very strong & decisive \\
         & mention & & & & \\
\bottomrule
\end{tabular}
\end{table}

\textbf{Proposal.} Replace the bare threshold with a minimum-evidence requirement
expressed on this ladder, and require the contract to name which rung constitutes
promotion. The value is contract-specific and cannot be defaulted.

\section{C3 and C4 --- Governance and operational constants}

\subsection{Staleness, freshness, epoch: reclassify, do not cite}

No literature was found for $3600$ (staleness), $600$ (freshness time-to-live),
or $86400$ (improvement epoch), and I do not believe any exists. These are not
scientific claims. A staleness threshold states how long an operator will wait
before treating an open record as abandoned; it is a property of the deployment's
cadence, not of nature.

The correct action is therefore \textbf{reclassification, not citation}: these
belong in the governance and operational classes, and the disclosure should say so
rather than list them beside a verdict confidence.

Two of them do carry a defect, but it is a duplication defect rather than a
missing one. The $600$ freshness value is written seven times across four
packages, so setting it in one place does not set it in the others. Log
retention is 100 files in four packages and 30 in one, meaning the audit trail
that gates all writes has different retention in one server than in the others.
Both are fixed by a single shared constants module and a stated reason, and
neither requires a citation.

\subsection{Operational remainders}

The observation grace period, the stalled-trial margin, and the $1.0$-second
margin that absorbs timestamp rounding are likewise ungrounded, and should be
recorded with a reason string rather than a reference: the last is a fact about
the clock, not about science. The default bind address is a security-posture
decision and belongs in an operator knob.

\section{Summary}
\label{sec:results}

\begin{table}[h]
\centering
\caption{All fifteen decision constants, with sourcing status and proposed
action.}
\label{tab:summary}
\small
\setlength{\tabcolsep}{3.5pt}
\begin{tabular}{@{}p{0.150\textwidth}p{0.068\textwidth}p{0.135\textwidth}p{0.527\textwidth}@{}}
\toprule
\textbf{Constant} & \textbf{Value} & \textbf{Status} & \textbf{Proposed action} \\
\midrule
prior confidence ceiling & 0.3 & \code{IN-RANGE} & keep; cite source; make prior deployable \\
verdict confidence & 0.85 & \code{UNGROUNDED} & replace with posterior; $0.624$ for prior .3, $p{=}0.05$ \\
meta-decision confidence & 0.6 & \code{UNGROUNDED} & delete; compute from evidence count \\
claim importance & 0.5 & \code{UNGROUNDED} & delete --- carries no information \\
edge weight & 1 & \code{UNGROUNDED} & rename to declare unranked, or drop \\
zero-divisor fallback & 0.0 & \code{VOID} & refuse; Fieller / Gleser--Hwang \\
seed count & 1 & \code{DERIVED} & $n \approx 9$ per arm (80\%); SESOI in contract \\
promotion threshold & $>1$ & \code{DERIVED} & minimum rung on the Bayes-factor ladder \\
staleness & 3600 & ungrounded & reclassify; deduplicate \\
freshness TTL & 600 & ungrounded & reclassify; deduplicate (7 copies) \\
improvement epoch & 86400 & ungrounded & reclassify \\
observation grace & 86400 & ungrounded & reclassify \\
stalled margin & 300 & ungrounded & reclassify \\
log retention & 100 / 30 & ungrounded & unify; document as audit policy \\
bind address & 0.0.0.0 & posture & operator knob \\
\bottomrule
\end{tabular}
\end{table}

\section{Acceptance criteria}

The proposal is implemented when:

\begin{enumerate}
\item \textbf{R1} --- every constant carries a source status from
\S\ref{sec:results} in both the code and the API schema.
\item \textbf{R2} --- no constant is defined in more than one location; a
continuous-integration check fails on the second definition.
\item \textbf{R3} --- the campaign closer refuses a zero divisor, matching the
tournament closer.
\item \textbf{R4} --- a tournament cannot open without an effect size of interest
and a target power in its promotion policy; the field stops being inert.
\item \textbf{R5} --- the frozen score row carries requested and achieved
sample size, the effect size of interest, and the sign and magnitude errors
where computable.
\item \textbf{R6} --- confidence is either computed from a stated prior and
likelihood, or is an enumeration.
\item \textbf{R7} --- constants with status \code{UNGROUNDED} may not be cited as
grounds for a promotion decision. This is the rule that makes the table above
binding rather than decorative.
\item \textbf{R8} --- a generated disclosure table --- value, file, line, class,
status, citation, last-changed commit --- is checked in and regenerated in CI.
\end{enumerate}

\section{What this proposal does not resolve}

\textbf{No literature fixes a staleness duration.} Any citation offered for the
value $3600$ is mislabelling a preference as a result.

\textbf{The calibration question is open, not solved.} \S\ref{sec:prior} specifies
how to \emph{compute} a defensible confidence; it does not say what this system's
prior ought to be. That is a domain judgement about agent-generated claims and
belongs to the maintainers.

\textbf{The power figures are model-based.} They assume roughly normal
per-observation noise. Whether they transfer to the heavy-tailed metric
distributions typical of deep-learning benchmarks is not established, and the
worked example in \S\ref{sec:n} inherits that assumption. The appropriate response
is simulation from the system's own observed trial distributions, not a different
textbook formula.

\textbf{One prior sub-test failed and is reported as failed.} A measurement of how
often a null arm mean lands on exactly zero returned $0.00\%$ and is
uninformative by construction, because continuous draws do not average to exactly
zero. The degenerate case arises from bounded or discrete metrics. That sub-test
has not been rebuilt.

\begin{thebibliography}{9}

\bibitem{nap2019}
National Academies of Sciences, Engineering, and Medicine.
\emph{Reproducibility and Replicability in Science}.
National Academies Press, 2019. Appendix D, Table D-1.

\bibitem{fieller1954}
E.~C. Fieller.
Some problems in interval estimation.
\emph{Journal of the Royal Statistical Society, Series B}, 16:174--185, 1954.

\bibitem{gleser1987}
L.~J. Gleser and J.~D. Hwang.
The application of Fieller's theorem for confidence
intervals for a ratio.
\emph{Biometrika}, 1987. Not obtained; known only through the restatement in
\cite{axioms2025}, which was read in full on the publisher's site.

\bibitem{axioms2025}
Re-examining confidence intervals for ratios of parameters.
\emph{Axioms}, 13(3):37, 2025. Read in full on the publisher's site; the PDF
could not be archived (the publisher's asset path resolved to a different
article in the same volume).

\bibitem{fiellerguide}
V.~H. Franz.
Ratios: a short guide to confidence limits and proper use.
Technical report, Justus-Liebig-Universit\"at Giessen, 2007. arXiv:0710.2024.

\bibitem{morey2016}
R.~D. Morey, R.~Hoekstra, J.~N. Rouder, M.~D. Lee, and E.-J. Wagenmakers.
The fallacy of placing confidence in confidence intervals.
\emph{Psychonomic Bulletin \& Review}, 23(1):103--123, 2016.

\bibitem{gelman2014}
A.~Gelman and J.~B. Carlin.
Beyond power calculations: assessing Type~S (sign) and Type~M (magnitude)
errors.
\emph{Perspectives on Psychological Science}, 9(6):641--651, 2014.

\bibitem{lakens2018sesoi}
D.~Lakens, A.~M. Scheel, and P.~M. Isager.
Equivalence testing for psychological research: a tutorial.
\emph{Advances in Methods and Practices in Psychological Science}, 1(2):259--269,
2018. Read in full on the publisher's site; the PDF is paywalled.

\bibitem{albers2018}
C.~J. Albers and D.~Lakens.
When power analyses based on pilot data are biased.
\emph{Journal of Experimental Social Psychology}, 2018.

\bibitem{kass1995}
R.~E. Kass and A.~E. Raftery.
Bayes factors.
\emph{Journal of the American\bs Statistical Association}, 90(430):773--795, 1995.

\bibitem{gneiting2007}
T.~Gneiting and A.~E. Raftery.
Strictly proper scoring rules, prediction, and estimation.
\emph{Journal of the American Statistical Association}, 102(477):359--378, 2007.

\bibitem{akker2023}
R.~van den Akker et al.
Preregistration in practice: a comparison of preregistered and
non-preregistered studies in psychology.
\emph{Behavior Research Methods}, 2023.

\bibitem{bakker2020}
B.~J. Bakker et al.
Recommendations in pre-registrations and internal review board proposals
promote formal power analyses but do not increase sample size.
\emph{PLOS ONE}, 2020.

\end{thebibliography}

\end{document}
